\subsection{Recurrent Neural Networks}
\label{sec:recurrent-neural-networks}
\slidesrc{Language models \& RNNs: 51--84}

\subsubsection{The idea: reuse the same weights at every step}
\slidesrc{Language models \& RNNs: 52--56}

A \textbf{recurrent neural network} (RNN) reads a text \emph{one word at a time}, the
way we do, and keeps a \textbf{hidden state} $\vect h^{(t)}$: a vector that summarises
everything read so far. At each step it combines the previous summary with the new
word to produce the new summary. Crucially, \textbf{the same weights are applied
repeatedly} at every step, so the network has the same size whether the text has 4
words or 4000: inputs can be of \textbf{any length}.

\begin{center}
\begin{tikzpicture}[>=Stealth,
    neu/.style={circle, draw=cyan!70!black, fill=cyan!35, minimum size=5pt, inner sep=0pt},
    hid/.style={draw=cyan!70!black, rounded corners, thick, minimum width=0.5cm,
                minimum height=1.45cm}]
  \foreach \w [count=\i] in {the, quick, brown, fox} {
    \pgfmathsetmacro{\x}{2.3*\i}
    \node[hid] (h\i) at (\x,0) {};
    \foreach \k in {1,...,4} \node[neu] at (\x,0.72-0.29*\k) {};
    \node[font=\scriptsize, anchor=south west, inner sep=1pt] at (h\i.north east) {$\vect h^{(\i)}$};
    \node[font=\small] (x\i) at (\x,-1.7) {\w};
    \node[font=\scriptsize, nbgray] at (\x,-2.05) {$\vect x^{(\i)}$};
    \node[font=\small] (y\i) at (\x,1.7) {$\hat{\vect y}^{(\i)}$};
    \draw[->, thick, nbgray] (x\i) -- (h\i);
    \draw[->, thick, nbgray] (h\i) -- (y\i);
  }
  \node[hid, draw=nbgray, minimum height=1.45cm] (h0) at (0,0) {};
  \node[font=\scriptsize, anchor=south west, inner sep=1pt] at (h0.north east) {$\vect h^{(0)}$};
  \foreach \i/\j in {0/1,1/2,2/3,3/4} {
    \draw[->, thick, nbgray] (h\i) -- node[draw=nbred, fill=white, font=\scriptsize,
        inner sep=2pt] (w\j) {$W_h$} (h\j);
  }
  \draw[->, thick, nbgray] (h4) -- ++(1.0,0) node[right] {\ldots};
  \node[font=\footnotesize, align=right, anchor=east] at (-0.6,1.7) {outputs};
  \node[font=\footnotesize, align=right, anchor=east] at (-0.6,0) {hidden\\states};
  \node[font=\footnotesize, align=right, anchor=east] at (-0.6,-1.85) {inputs, of\\any length};
  \node[draw=nbred, rounded corners, fill=nbred!5, font=\footnotesize, align=center] at (5.75,-2.9)
       {all the red boxes are the \emph{same} weights $W_h$,\\applied repeatedly to the whole sequence};
\end{tikzpicture}
\captionof{figure}{An RNN ``unrolled'' over four time steps. There is only one matrix
$W_h$ (and one $W_e$, one $U$): the copies are the same weights used again at every
step.}
\end{center}

\begin{nbdef}{RNN language model}
At each time step $t = 1, 2, \ldots$
\begin{align*}
  \vect e^{(t)} &= E \vect x^{(t)}
      && \text{word embedding of the one-hot word } \vect x^{(t)} \in \mathbb R^{|V|}\\
  \vect h^{(t)} &= \sigma\bigl(W_h \vect h^{(t-1)} + W_e \vect e^{(t)}\bigr)
      && \text{new hidden state} \\
  \hat{\vect y}^{(t)} &= \softmax\bigl(U \vect h^{(t)}\bigr) \in \mathbb R^{|V|}
      && \text{distribution of the next word } \vect x^{(t+1)}
\end{align*}
$\vect h^{(0)}$ is an initial hidden state (often a vector of zeros). The parameters are
$E$, $W_h$, $W_e$, $U$, shared by all time steps.
\end{nbdef}

Read the hidden-state equation slowly: $W_h \vect h^{(t-1)}$ is ``what I remember'',
$W_e \vect e^{(t)}$ is ``what I just read''; they are added and squashed by $\sigma$ into
the new memory $\vect h^{(t)}$. Since $\vect h^{(t-1)}$ itself depended on
$\vect h^{(t-2)}$ and on $\vect x^{(t-1)}$, and so on, $\vect h^{(t)}$ depends, in
principle, on \emph{all} the words $\vect x^{(1)}, \ldots, \vect x^{(t)}$. There is no
window any more.

\begin{nbintuition}[an RNN is a reader taking notes]
\notinslides
Imagine reading a novel with a notebook of fixed size. After each word you update your
notes (same rules every time), and from the notes alone you guess the next word. The
notebook never grows, whatever the length of the book, so you must decide what to keep
and what to overwrite. This is both the strength of RNNs (any length, constant size)
and, as we will see, their weakness (old information gets overwritten).
\end{nbintuition}

\needspace{12\baselineskip}
\subsubsection{Pros and cons}
\slidesrc{Language models \& RNNs: 57}

\begin{center}
\small
\begin{tabularx}{\linewidth}{@{}X X@{}}
\toprule
\textbf{Pros} & \textbf{Cons} \\
\midrule
not restricted to fixed-length inputs &
  computation is slow: step $t$ needs $\vect h^{(t-1)}$, so the steps cannot be computed
  in parallel \\
can use information from many time steps ago (in theory) &
  in practice long-range information is hard to capture
  (Section~\ref{sec:vanishing-exploding-gradients}) \\
input words are processed symmetrically: the same weights at every position
(\textbf{weight sharing}), so what is learnt about a word in one position applies
everywhere & \\
\bottomrule
\end{tabularx}
\end{center}
Compare with the fixed-window model (Section~\ref{sec:neural-language-models}): the
first and third pros solve exactly its ``fixed window'' and ``no symmetry'' problems,
and the model size no longer depends on the length of the context.

\subsubsection{Training an RNN language model}
\slidesrc{Language models \& RNNs: 58--67}

\begin{enumerate}
  \item Gather a large corpus of text $\vect x^{(1)}, \ldots, \vect x^{(T)}$.
  \item Feed it to the RNN one word at a time, computing $\hat{\vect y}^{(t)}$ for every
        $t = 1, \ldots, T$.
  \item At each step compute the loss: the cross-entropy between the predicted
        distribution $\hat{\vect y}^{(t)}$ and the true next word
        $\vect y^{(t)} = \vect x^{(t+1)}$ (one-hot):
        \[
          J^{(t)} = \mathrm{CE}(\vect y^{(t)}, \hat{\vect y}^{(t)})
          = -\sum_{w \in V} \underbrace{y^{(t)}_w}_{\text{one-hot}} \log \hat y^{(t)}_w
          = -\log \hat y^{(t)}_{\vect x^{(t+1)}} ,
        \]
        the negative log-probability of the correct next word.
  \item Average the loss over the whole dataset (the last word has no next word to
        predict, hence $T-1$ terms):
        \[
          J = \frac{1}{T-1} \sum_{t=1}^{T-1} J^{(t)}
            = \frac{1}{T-1} \sum_{t=1}^{T-1} -\log \hat y^{(t)}_{\vect x^{(t+1)}} .
        \]
\end{enumerate}

\begin{center}
\begin{tikzpicture}[>=Stealth,
    neu/.style={circle, draw=#1!70!black, fill=#1!35, minimum size=4pt, inner sep=0pt},
    box/.style={draw=#1!70!black, rounded corners, minimum width=0.42cm, minimum height=1.15cm}]
  \foreach \w/\nx [count=\i] in {the/quick, quick/brown, brown/fox, fox/jumped} {
    \pgfmathsetmacro{\x}{2.2*\i}
    \node[box=magenta] (e\i) at (\x,0) {};
    \foreach \k in {1,...,4} \node[neu=magenta] at (\x,0.57-0.23*\k) {};
    \node[box=cyan] (h\i) at (\x,1.9) {};
    \foreach \k in {1,...,4} \node[neu=cyan] at (\x,2.47-0.23*\k) {};
    \node[font=\small] (x\i) at (\x,-1.3) {\w};
    \node[font=\scriptsize] (y\i) at (\x,3.25) {$\hat{\vect y}^{(\i)}$};
    \node[font=\scriptsize, nbred] (J\i) at (\x,4.05) {$J^{(\i)}$};
    \draw[->, nbgray] (x\i) -- node[right, font=\tiny] {$E$} (e\i);
    \draw[->, nbgray] (e\i) -- node[right, font=\tiny] {$W_e$} (h\i);
    \draw[->, nbgray] (h\i) -- node[right, font=\tiny] {$U$} (y\i);
    \draw[->, nbred] (y\i) -- (J\i);
    \node[font=\tiny, nbred, align=center] at (\x,4.55) {$-\log \hat y^{(\i)}_{\textit{\nx}}$};
  }
  \node[box=cyan, draw=nbgray] (h0) at (0,1.9) {};
  \node[font=\scriptsize] at (0,2.75) {$\vect h^{(0)}$};
  \foreach \i/\j in {0/1,1/2,2/3,3/4} \draw[->, nbgray] (h\i) -- node[above, font=\tiny] {$W_h$} (h\j);
  \draw[->, nbgray] (h4) -- ++(0.8,0) node[right] {\ldots};
  \node[font=\small] at (11.0,-1.3) {jumped};
  \node[font=\scriptsize, nbred, anchor=east] at (-0.3,4.05) {loss at each step};
  \node[font=\scriptsize, nbred, anchor=east] at (-0.3,3.25) {predicted distributions};
  \node[font=\scriptsize, nbred, anchor=east] at (-0.3,-1.3) {corpus};
  \node[font=\small, nbred] at (5.5,5.2)
       {$J = \frac14\bigl(J^{(1)} + J^{(2)} + J^{(3)} + J^{(4)}\bigr)$};
\end{tikzpicture}
\captionof{figure}{Training on ``the quick brown fox jumped''. At step 1 the network has
read \emph{the} and is penalised by $-\log$ of the probability it gave to
\emph{quick}; at step 2 by the probability it gave to \emph{brown}; and so on. The total
loss is the average.}
\end{center}

Note that the input at every step is the \emph{true} word from the corpus, not the
word the model predicted: during training the model is always shown the correct text
and asked to predict one step ahead.

\paragraph{Training in practice: stochastic gradient descent.}
A typical corpus has millions of words. Computing the loss
$J(\theta) = \frac{1}{T-1}\sum_{t} J^{(t)}(\theta)$ and its gradient over all of it at
once requires too much memory and time. So:
\begin{itemize}[nosep]
  \item divide the corpus into smaller \textbf{chunks} (sentences, documents);
  \item \textbf{Stochastic Gradient Descent} (SGD) approximates the full gradient with
        the average gradient over a small \textbf{batch} of $S$ chunks:
        \[
          \nabla_\theta J(\theta) \approx \frac1S \sum_{j=1}^{S} \nabla_\theta J^{(j)}(\theta) ;
        \]
  \item loss and gradients are computed for one batch of sentences at a time, then the
        weights are updated ($\theta \leftarrow \theta - \alpha \nabla_\theta J$) and the
        next batch is taken.
\end{itemize}
The approximation is noisy (a different batch gives a slightly different gradient),
but on average it points in the right direction, and each step is cheap. That noise is
the ``stochastic'' in the name.

\subsubsection{Backpropagation through time}
\slidesrc{Language models \& RNNs: 68--76}

To run gradient descent we need the gradient of the loss with respect to every
parameter. The tricky one is $W_h$, because it is used at \emph{every} step. To see the
structure clearly, the slides use an even simpler model with no input and no
non-linearity:
\[
  \vect h^{(t)} = W_h \vect h^{(t-1)}, \qquad
  J^{(t)}(\theta) = \mathrm{CE}(\vect h^{(t)}, \vect y^{*(t)}) .
\]

\begin{nbfact}{Backpropagation through time (BPTT)}
The gradient with respect to a repeated weight is the \textbf{sum of the gradients for
each time it appears}:
\[
  \frac{\partial J^{(t)}}{\partial W_h} =
  \sum_{i=1}^{t} \left. \frac{\partial J^{(t)}}{\partial W_h} \right|_{(i)},
\]
where $\bigl.\frac{\partial J^{(t)}}{\partial W_h}\bigr|_{(i)}$ is the gradient obtained by
treating only the copy of $W_h$ used at step $i$ as variable.
\end{nbfact}

\begin{nbintuition}[why a sum?]
\notinslides
Think of $W_h$ as a knob that turns $t$ dials at once (one per step). Turning it a
little changes the output through each dial; for small changes the effects simply add
up. So the total sensitivity is the sum of the sensitivities through each use.
\end{nbintuition}

\paragraph{The derivation.}
Start with the \textbf{chain rule}: the loss depends on $W_h$ only through
$\vect h^{(t)}$, so
\[
  \frac{\partial J^{(t)}}{\partial W_h} =
  \frac{\partial J^{(t)}}{\partial \vect h^{(t)}}\,
  \frac{\partial \vect h^{(t)}}{\partial W_h}.
\]
The first factor is easy (it only involves the output). The second is the interesting
one: $\vect h^{(t)} = W_h \vect h^{(t-1)}$ depends on $W_h$ in two ways, \emph{directly}
(it is multiplied by $W_h$) and \emph{indirectly} (through $\vect h^{(t-1)}$, which was
itself computed with $W_h$). Write it as $\vect h^{(t)}(W_h, \vect h^{(t-1)})$.

To keep the algebra readable we first do it with \textbf{numbers instead of matrices}
(a hidden state of size 1, so $h^{(t)} = w\, h^{(t-1)}$). The derivative of a product is
``derivative of the first times the second plus the first times the derivative of the
second'' (product rule), so:
\begin{align*}
  \frac{\dd h^{(t)}}{\dd w}
  &= \underbrace{h^{(t-1)}}_{\text{direct}}
     + \underbrace{w\, \frac{\dd h^{(t-1)}}{\dd w}}_{\text{indirect}}
     \reason{product rule on $w\,h^{(t-1)}$} \\
  &= h^{(t-1)} + w \Bigl( h^{(t-2)} + w\, \frac{\dd h^{(t-2)}}{\dd w} \Bigr)
     \reason{same rule for $h^{(t-1)}$} \\
  &= h^{(t-1)} + w\, h^{(t-2)} + w^2\, \frac{\dd h^{(t-2)}}{\dd w}
     \reason{expand} \\
  &= \cdots = h^{(t-1)} + w\, h^{(t-2)} + w^2 h^{(t-3)} + \cdots + w^{t-1} h^{(0)}
     \reason{$\dd h^{(0)}/\dd w = 0$} \\
  &= \sum_{i=0}^{t-1} w^{\,t-1-i}\, h^{(i)} .
\end{align*}
Each term is one ``use'' of $w$: the term with $h^{(i)}$ comes from the copy of $w$
that multiplies $h^{(i)}$ at step $i+1$, and the factor $w^{t-1-i}$ carries its effect
through the remaining $t-1-i$ steps up to $h^{(t)}$. This is the BPTT sum.

\begin{nbexample}{checking the formula}
\notinslides
Take $w = 0.7$, $h^{(0)} = 2$, $t = 5$. Since $h^{(t)} = w^t h^{(0)}$, directly
$\frac{\dd h^{(5)}}{\dd w} = 5 w^4 h^{(0)} = 5 \cdot 0.2401 \cdot 2 = 2.401$. With the
formula, $h^{(i)} = 2 \cdot 0.7^i$ and every term equals $w^{4-i} \cdot 2 w^i = 2w^4$:
\[
  \sum_{i=0}^{4} w^{4-i} h^{(i)} = 5 \cdot 2 \cdot 0.7^4 = 2.401 . \quad\checkmark
\]
\end{nbexample}

With matrices the computation is the same, and the slides write the result as
\[
  \frac{\partial \vect h^{(t)}}{\partial W_h}
  = \sum_{i=0}^{t-1} \Bigl( \prod_{j=i+1}^{t-1} W_h \Bigr) \vect h^{(i)}
  = \sum_{i=0}^{t-1} W_h^{\,t-1-i}\, \vect h^{(i)} .
\]

\begin{nbwarning}[what the matrix formula means]
The derivative of a vector with respect to a matrix is a three-index object, so the
slide's formula (and its intermediate step ``$\vect h^{(t-1)} I$'') is a shorthand. The
precise statement is: if $W_h$ is changed by a small matrix $\Delta W$, then
$\vect h^{(t)}$ changes by
\[
  \Delta \vect h^{(t)} \approx \sum_{i=0}^{t-1} W_h^{\,t-1-i}\, \Delta W\, \vect h^{(i)} .
\]
The structure is exactly the scalar one: one term per use of $W_h$, each carried to step
$t$ by a power of $W_h$.
\end{nbwarning}

\begin{itemize}[nosep]
  \item BPTT is \textbf{equivalent to ordinary backpropagation through the whole
        unrolled network}: once the RNN is drawn as a long chain of copies, it is just a
        (very deep) feed-forward network whose copies share the same weights.
  \item In practice the backward pass is \textbf{truncated} at about 20 steps, for
        efficiency: terms older than that are ignored (\textbf{truncated BPTT}).
  \item For more complex RNNs (with inputs and non-linearities) the formula is very
        similar.
\end{itemize}

\subsubsection{Generating text with an RNN}
\slidesrc{Language models \& RNNs: 77}

To generate text at test time, \textbf{repeatedly sample} from the network,
\textbf{feeding each output back as the next input}:
\begin{enumerate}[nosep]
  \item feed the start token \texttt{<s>}; the network outputs $\hat{\vect y}^{(1)}$;
  \item sample a word from $\hat{\vect y}^{(1)}$, say \emph{the};
  \item feed \emph{the} as the next input; the network updates its hidden state and
        outputs $\hat{\vect y}^{(2)}$; sample \emph{quick};
  \item feed \emph{quick}, sample \emph{brown}, and so on, until \texttt{</s>} is
        sampled or a length limit is reached.
\end{enumerate}
Unlike an n-gram model, the hidden state carries information from \emph{all} the words
generated so far, not only the last $n-1$. This is also, in essence, how chatbots
write their answers, one token at a time.

\subsubsection{Evaluating language models: perplexity}
\slidesrc{Language models \& RNNs: 78--84}

How do we measure whether a language model is good? We take a \textbf{test sentence}
$X_{\text{test}} = \vect x^{(1)}, \ldots, \vect x^{(T)}$ that was \textbf{not seen
during training} (otherwise a model that memorised the corpus would look perfect) and
check how much probability the model gives it.

\begin{nbdef}{Perplexity}
The \textbf{perplexity} of a sentence is its inverse probability, normalised by its
length $T$:
\[
  \text{perplexity}(X_{\text{test}}) = p(\vect x^{(1)}, \ldots, \vect x^{(T)})^{-1/T}
  = \sqrt[T]{\frac{1}{p(\vect x^{(1)}, \ldots, \vect x^{(T)})}} .
\]
\end{nbdef}
To compute it:
\begin{enumerate}[nosep]
  \item compute the probability of the sentence with the tested model and the chain
        rule, $p(\vect x^{(1)}, \ldots, \vect x^{(T)}) = \prod_{t=1}^{T}
        p(\vect x^{(t)} \mid \vect x^{(1)}, \ldots, \vect x^{(t-1)})$;
  \item raise it to the power $-1/T$.
\end{enumerate}
Why both pieces?
\begin{itemize}[nosep]
  \item \textbf{Inverse}: a good model gives the test sentence a \emph{high}
        probability, hence a \emph{low} perplexity. Lower is better.
  \item \textbf{$T$-th root}: the probability of a sentence shrinks with every word, so
        long sentences would always look worse. The $T$-th root is the
        \emph{geometric mean} per word, which makes sentences of different lengths
        comparable.
\end{itemize}

\begin{nbintuition}[perplexity as a number of choices]
\notinslides
A model that spreads its probability uniformly over $k$ words at every step gives the
sentence probability $(1/k)^T$, so its perplexity is
$\bigl((1/k)^T\bigr)^{-1/T} = k$. Perplexity is therefore the \emph{effective number of
equally likely words} the model is hesitating between at each step. A perplexity of 50
means ``as confused as if choosing among 50 words''. A model that knows nothing
(uniform over the whole vocabulary) has perplexity $|V|$.
\end{nbintuition}

\paragraph{Perplexity in practice and its link with the loss.}
\begin{itemize}[nosep]
  \item Computing $\prod_{t=1}^{T} p(\vect x^{(t)} \mid \vect x^{(1)}, \ldots,
        \vect x^{(t-1)})$ for long sequences is \textbf{numerically unstable}: the
        product of thousands of numbers smaller than 1 becomes too small for the
        computer to represent (it rounds to 0). In practice we include begin and end
        markers \texttt{<s>}, \texttt{</s>} and \textbf{truncate the conditioning} at
        sentence or document boundaries, and we work with logarithms.
  \item \textbf{Connection to the cross-entropy loss.} For an RNN, the probability of
        each word is read from the output of the previous step:
        $p(\vect x^{(t+1)} \mid \vect x^{(1)}, \ldots, \vect x^{(t)}) =
        \hat y^{(t)}_{\vect x^{(t+1)}}$. Using $a = e^{\log a}$ and
        $\log \prod = \sum \log$:
        \begin{align*}
          \text{perplexity}
          &= \Bigl[\prod_{t=1}^{T} \hat y^{(t)}_{\vect x^{(t+1)}}\Bigr]^{-1/T}
           = \exp\Bigl( -\frac1T \log \prod_{t=1}^{T} \hat y^{(t)}_{\vect x^{(t+1)}} \Bigr)\\
          &= \exp\Bigl( \frac1T \sum_{t=1}^{T}
             \underbrace{-\log \hat y^{(t)}_{\vect x^{(t+1)}}}_{J^{(t)}} \Bigr)
           = \exp(\text{average cross-entropy loss}) .
        \end{align*}
        Training an LM by minimising the cross-entropy loss is the same as minimising
        perplexity.
  \item \textbf{Lower perplexity} means the model is \textbf{less ``surprised''} by the
        text it is given.
\end{itemize}

\begin{nbwarning}[indices on the slide]
The slide writes $p(\vect x^{(t)} \mid \vect x^{(1)}, \ldots, \vect x^{(t-1)}) =
\hat y^{(t)}_{\vect x^{(t+1)}}$, mixing two conventions, and sums over $i$ while the
terms depend on $t$. The consistent version is the one above: $\hat{\vect y}^{(t)}$ is
computed after reading $\vect x^{(t)}$ and predicts $\vect x^{(t+1)}$, so the $t$-th
factor is $\hat y^{(t)}_{\vect x^{(t+1)}}$. Here the sequence starts with
$\vect x^{(1)} = \texttt{<s>}$ (which is given, not predicted) and the $T$ predicted
tokens are $\vect x^{(2)}, \ldots, \vect x^{(T+1)}$, the last one being \texttt{</s>}.
\end{nbwarning}

\begin{nbexample}{perplexity of two models}
\slidesrc{Language models \& RNNs: 82--84}
The vocabulary is the digits 0 to 9, $|V| = 10$.
\begin{itemize}[nosep]
  \item LM1 outputs any digit with probability $\frac{1}{10}$;
  \item LM2 outputs 0 with probability $\frac{91}{100}$ and each other digit with
        probability $\frac{1}{100}$ (check: $0.91 + 9 \cdot 0.01 = 1$).
\end{itemize}
Which one has the lowest perplexity on the sequence ``0 0 0 3 0'' ($T = 5$)?
\begin{align*}
  \text{LM1:}\quad \text{perplexity} &= \Bigl[\bigl(\tfrac{1}{10}\bigr)^5\Bigr]^{-1/5}
     = \bigl(\tfrac{1}{10}\bigr)^{-1} = 10 \\
  \text{LM2:}\quad \text{perplexity} &= \Bigl[\bigl(\tfrac{91}{100}\bigr)^4 \cdot
     \tfrac{1}{100}\Bigr]^{-1/5} = (0.006857)^{-1/5} \approx 2.709
\end{align*}
LM2 wins: it expects zeros, the sequence is mostly zeros, so it is much less surprised.
LM1 is uniform, so its perplexity is $|V| = 10$, as predicted by the intuition box.
Note that perplexity depends on the test text: on ``3 3 3 3 3'' LM2 would get
$\bigl((1/100)^5\bigr)^{-1/5} = 100$, ten times worse than LM1.
\end{nbexample}

\begin{nbexercise}{2}{perplexity of the toy bigram model}
\notinslides
Compute the perplexity of the bigram model of Section~\ref{sec:ngram-models} (toy
corpus, no smoothing) on the sentence ``I like dogs'', counting the end marker as a
predicted word.
\end{nbexercise}
\begin{nbsolution}
The model predicts $T = 4$ tokens: \emph{I}, \emph{like}, \emph{dogs}, \texttt{</s>},
with probabilities $\frac23$, $1$, $\frac13$, $1$. The sentence probability is
$\frac29$, so
\[
  \text{perplexity} = \bigl(\tfrac29\bigr)^{-1/4} = 4.5^{1/4} \approx 1.46 .
\]
Very low: the toy model has almost no choices to make (after \emph{I} it is certain of
\emph{like}). Lesson: perplexity is only meaningful when comparing models on the
same, realistic, unseen test text; a tiny model evaluated on its own training sentences
looks deceptively good.
\end{nbsolution}

\begin{nbsummary}[RNN language models]
\begin{itemize}[nosep]
  \item $\vect h^{(t)} = \sigma(W_h \vect h^{(t-1)} + W_e E \vect x^{(t)})$,
        $\hat{\vect y}^{(t)} = \softmax(U \vect h^{(t)})$; same weights at every step.
  \item Loss: $J = \frac{1}{T-1}\sum_t -\log \hat y^{(t)}_{\vect x^{(t+1)}}$, minimised
        with SGD on batches.
  \item BPTT: gradient w.r.t.\ a shared weight = sum over its uses,
        $\frac{\partial \vect h^{(t)}}{\partial W_h} \sim \sum_{i} W_h^{\,t-1-i}
        \vect h^{(i)}$; truncated at $\sim$20 steps.
  \item Generation: sample, feed the sample back as input, repeat.
  \item Perplexity $= p(X)^{-1/T} = \exp(\text{average CE loss})$; lower is better;
        uniform model $\Rightarrow$ $|V|$.
\end{itemize}
\end{nbsummary}
